% --- ARCHIVO: Anexos_Conceptos/anexo_cap6.tex ---

\chapter{Conceptos sobre comunicación de crisis del Capítulo 6}
\label{anexo:conceptos_cap6}

\section{Crisis communication failure}
\label{anexo:crisis_comm_failure}
En gestión de emergencias, un \textit{crisis communication failure} describe el fallo institucional al no ocupar de forma oportuna el espacio informativo con mensajes verificables tras un incidente grave. Según el \textit{Chaos Communication Model}, si la institución responsable no emite un relato claro durante la ventana crítica inicial (1-6 horas), el vacío discursivo es ocupado por narrativas alternativas o no verificadas que, al consolidarse, resultan casi imposibles de revertir mediante correcciones técnicas posteriores.

\section{Infodemia}
\label{anexo:infodemia}
Término popularizado por la OMS para describir la sobrepoblación del espacio informativo con contenidos no verificados, erróneos o falsos que se propagan rápidamente en situaciones de crisis. En incidentes de infraestructuras críticas, la infodemia se acelera por la asimetría entre la inmediatez para crear rumores y la lentitud para verificar datos técnicos, agravándose cuando la propia naturaleza del evento (ej. un apagón) priva a la población del acceso normalizado a fuentes de contraste fiables.

\section{Encuadre mediático (Framing) y Agenda-shifting}
\label{anexo:framing}
El \textit{framing} es el proceso por el cual los medios seleccionan y enfatizan ciertos elementos de un hecho para proponer una interpretación causal concreta, influyendo en a quién atribuye la responsabilidad el público. Relacionado con esto, el \textit{agenda-shifting} ocurre cuando un evento disruptivo es instrumentalizado mediáticamente para desplazar la atención y reabrir debates políticos o estructurales preexistentes (como el modelo energético o el cierre nuclear), que no fueron causados por el incidente en sí.

\section{Vacuum filling (Relleno del vacío informativo)}
\label{anexo:vacuum_filling}
El \textit{vacuum filling} es el proceso estructural e inevitable mediante el cual la incertidumbre colectiva ante un desastre genera una demanda de respuestas que, si no es satisfecha por las instituciones oficiales, es cubierta espontáneamente por fuentes no autorizadas. Según la teoría SCCT (Coombs), las explicaciones causales que llenan este vacío en las primeras horas tienden a arraigar profundamente en la percepción pública, condicionando la narrativa de la crisis a largo plazo.

\section{Emergent norm theory}
\label{anexo:emergent_norm}
La teoría de las normas emergentes (Turner y Killian) sostiene que, frente a las visiones del pánico masivo, los grupos en situaciones de disrupción desarrollan espontáneamente nuevas reglas de comportamiento social adaptativo. En el ecosistema digital durante un apagón, esto se traduce en la creación rápida de redes improvisadas de asistencia mutua, intercambio de información útil de proximidad y el uso del humor o la ironía como mecanismo colectivo para regular el estrés emocional.

\section{Outrage communication (Comunicación de indignación)}
\label{anexo:outrage_comm}
Basado en el modelo de Sandman (\textit{Risk = Hazard + Outrage}), este concepto indica que la percepción pública de un riesgo depende más de factores emocionales (indignación, percepción de negligencia) que de la evaluación técnica del peligro real. En redes sociales, los algoritmos incentivan este patrón al viralizar mensajes de alta carga emocional y polarización, desincentivando la divulgación técnica aséptica, lo que desplaza el debate de las causas físicas a la atribución de culpas políticas.